% !PRINTED 503
\weberpage{503}

Since, as we have already seen above, two subsidiary groups of $Q$ are either wholly identical, or contain no common element, we can further state the theorem:

2. \gesperrt{If $\pi$ is an arbitrary permutation from $P$, then the system of subsidiary groups}
\[
Q\pi,\ Q\pi_1\pi,\ Q\pi_2\pi\ldots Q\pi_{j-1}\pi
\]
\gesperrt{is distinguished from the system}
\[
Q,\ Q\pi_1,\ Q\pi_2\ldots Q\pi_{j-1}
\]
\gesperrt{only by the arrangement.}

We extend the designation subsidiary groups also to systems of the form $\pi Q$, and can decompose $P$ into such systems as well. From the decomposition (4) there follows in particular the second
\[
P = Q + \pi_1^{-1}Q + \pi_2^{-1}Q + \cdots \pi_{j-1}^{-1}Q.
\]
For if $\pi_1$ does not occur in $Q$, then $\pi_1^{-1}$ does not occur in it either, and we obtain a first subsidiary group $\pi_1^{-1}Q$. If then $\pi_2$ does not occur in $Q\pi_1$, then $\pi_2^{-1}$ does not occur in $\pi_1^{-1}Q$ either; for from $\pi_2^{-1} = \pi_1^{-1}\varkappa$ would follow $\pi_2 = \varkappa^{-1}\pi_1$, and so forth.

We here give the preference to the decomposition (4) on account of the significance which, as we shall see presently, it has for the algebraic problem.

We call a quantity $\psi$ of the field $N$ \gesperrt{belonging to the group $Q$} (when necessary with the addition ``within $P$''), if $\psi$ does not change when one of the operations of $Q$ is applied to it, but does change when any other operation from $P$ is applied. As one sees, this concept of belonging is somewhat extended as against the one fixed in \S.\,152, inasmuch as the permutations that lie outside $P$ do not come into consideration here at all. For the case that $P$ is the symmetric group, however, the two definitions coincide.

With this definition of belonging, the following theorems now hold generally.

3. \gesperrt{To every divisor $Q$ of $P$ there belong quantities in $N$, and every quantity in $N$ belongs to a determinate divisor $Q$ of $P$.}

The first part of the theorem is already proved in \S.\,152, 5. For if a function of the quantities $\alpha$ within the sym-
%CONT>>
